\documentclass[11pt]{article}
\usepackage[a4paper,margin=1in]{geometry}
\usepackage{amsmath,amssymb,amsthm}
\usepackage{hyperref}
\usepackage{needspace}
\newtheorem{theorem}{Theorem}
\newtheorem{lemma}[theorem]{Lemma}
\newtheorem{proposition}[theorem]{Proposition}
\newcommand{\FF}{\mathbb{F}}
\newcommand{\cF}{\mathcal{F}}
\title{Zero-Sum Quartet Incidence and Sharp Maximal-Independent-Set Witnesses on Eight Vertices}
\author{John Fairfax-Ball}
\date{Research draft, 16 October 2026 --- NOT FOR SUBMISSION}
\begin{document}
\maketitle

\begin{abstract}
We combine the four-quartet incidence pattern of an eight-point cap in
$\FF_3^3$ with maximal independent sets of connected graphs on eight
vertices. Assuming the classical affine equivalence classification of
nine-caps, we derive the complete distribution of the number of zero-sum
four-vertex maximal independent sets under a uniform cap labelling.
Its three-witness probability is at most $1/105$, with equality precisely
for the spider tree with leg lengths $(2,2,2,1)$. The finite $4\times4$
mixed-maximal-independent-set extremal component has a
Lean-verified bit-vector foundation; its three arbitrary Boolean-matrix semantic extremal transfers
have been strictly Lean-kernel-verified on research pull request \#3
at commit \href{https://github.com/jfairfaxball-348/Pan-Disciplinary-Invariance-Conjecture/commit/1e4cadd580b385e0d8e69e836f3211129e705839}{\texttt{1e4cadd580b3}}. This draft
makes no claim of literature priority or a complete end-to-end
formalisation.
\end{abstract}

\section{Introduction and scope}
Our central statistic is not the total number of maximal independent sets
of a graph. Instead it counts the maximal independent sets of cardinality
four that coincide with zero-sum quartets of an algebraic eight-cap.
We separate the algebraic cap input from the independent finite graph
argument and from the random-labelling probability model.
Relevant maximal-independent-set extremal work includes Liu
\cite{Liu1993}; our novelty relative to that literature remains to be
determined by source-level comparison.

\section{Classical cap input and quartet incidence}
Let $U\subset\FF_3^3\setminus\{0\}$ be an eight-point admissible cap, i.e.,
one with no zero-sum subsets of cardinality at most three. We use the
classical equivalence of nine-caps in $AG(3,3)$ from
Follett--Kalail--McMahon--Pelland--Won \cite[Proposition 2.2]{Follett2014}
as an explicitly cited external input.

\begin{lemma}[Quartet pattern; cap-classification input]
There are four four-element zero-sum subsets of $U$. Their unordered
incidence pattern is a four-cycle: two disjoint complementary pairs, and
every other pair overlaps in two points. Every triple of distinct quartets
contains a unique disjoint pair.
\end{lemma}
\begin{proof}
Augment $U$ by $0$ to obtain a nine-cap: lines through $0$ would
correspond to the forbidden opposite pairs in $U$. A canonical nine-cap is
$W_0=\{(x,y,x^2+y^2):(x,y)\in\FF_3^2\}$. Write $q(x)=x_1^2+x_2^2$.
The affine automorphisms
$\phi_t(x,z)=(x+t,z+2\langle x,t\rangle+q(t))$
preserve $W_0$ and act transitively on its nine points. Compose the
classical affine equivalence $U\cup\{0\}\to W_0$ with an appropriate
$\phi_t$ so that $0$ maps to $0$. The resulting equivalence is linear;
this is essential since translation need not preserve four-point zero sums.

In $W_0\setminus\{0\}$ there are four axis points at height $1$ and
four corner points at height $2$. A zero-sum quartet must have two of
each: if it contains $a$ axis points, its height sum is $8-a$, forcing
$a=2$. Two axis projections sum either to $0$ or to a vector with two
nonzero coordinates; two corner projections sum either to $0$ or to a
vector with one nonzero coordinate. Cancellation therefore requires
both projected pairs to be opposite. Let $A_1,A_2$ denote the lifted
axis pairs and $B_1,B_2$ the lifted corner pairs. The four quartets are
exactly $A_i\cup B_j$. Their intersections give the claimed pattern,
which transports under the linear equivalence above.
\end{proof}

\section{The exact mixed-witness invariant}
Fix a Boolean matrix $A=(A_{ij})_{i,j\in[4]}$, viewed as the edges of a
bipartite graph with four vertices on each side. Let
\[
e(A)=\#\{(i,j):A_{ij}=1\}.
\]
Let $\mathcal M(A)$ consist of pairs $(X,Y)$ of two-element subsets
of the left and right sides such that $X\cup Y$ is maximal independent.
Set
\[
b(A)=\#\{X:\exists Y,\ (X,Y)\in\mathcal M(A)\}.
\]
For $X\subseteq[4]$, let $N(X)$ be its opposite-side neighbourhood.

\begin{lemma}[Unique opposite-side partner]
For any mixed $2+2$ maximal independent set $X\cup Y$,
\[
Y=[4]\setminus N(X).
\]
Moreover $X$ is counted by $b(A)$ if and only if
\[
|N(X)|=2,\quad
\text{and}\quad
\forall i\notin X\ \exists j\in[4]\setminus N(X):A_{ij}=1.
\]
\end{lemma}
\begin{proof}
Independence makes $Y$ disjoint from $N(X)$; maximality forces
every vertex outside $Y$ on the right to lie in $N(X)$. The remaining
maximality conditions give the displayed criterion.
\end{proof}

\begin{theorem}[Sharp finite mixed-witness bounds]
For every Boolean $4\times4$ matrix $A$,
\[
 b(A)\ge4\ \Longrightarrow\ e(A)\le5,\qquad
 e(A)\ge8\ \Longrightarrow\ b(A)\le2.
\]
Furthermore,
\[
 e(A)=7,\quad b(A)=3
 \quad\Longleftrightarrow\quad
 A\text{ is a row- or column-centred spider matrix}.
\]
\end{theorem}
\begin{proof}[Proof outline]
Encode each row as four bits. Exhaustive \emph{kernel-checked}
finite-proposition proofs establish the three BitVec implications,
while independent symbolic left-pair/edge cardinality bridges show
the encoding gives exactly $b(A)$ and $e(A)$. The semantic transfer
statements passed the exact-SHA S003R8 hosted Lean workflow
38887318919, with only standard axioms. S003R9 constructs the associated simple graph on
$[4]\sqcup[4]$ and a bijection between its actual mixed maximal
independent vertex sets and the left-pair witnesses. These declarations passed the S003R9 strict kernel workflow
38889327092 at source revision \texttt{8199e6bc41dc}. The full connected
triple-count and probability theorems remain outside that certificate.
\end{proof}

\section{From complementary sets to the spider}
For a connected simple eight-vertex graph $G$, define
$\cF(G)$ to be its size-four maximal independent sets. Let $h(G)$
count unordered triples $\{S,T,R\}\subseteq\cF(G)$ for which
$S\cap T=\varnothing$ and $|R\cap S|=|R\cap T|=2$.

\begin{theorem}[Sharp graph extremiser]
For every connected simple eight-vertex graph,
\[
h(G)\le3,\qquad
h(G)=3\ \Longleftrightarrow\
G\cong S(2,2,2,1).
\]
\end{theorem}
\begin{proof}
If $h(G)>0$, the complementary disjoint independent sets $S,T$
form a balanced bipartition. By connectedness it is the unique
bipartition up to swap and both classes are maximal independent.
Every contributing third set has a $2+2$ split, so $h(G)=b(A)$ for
the associated $4\times4$ adjacency. Connectivity gives at least
seven edges; the finite theorem forces $h\le3$ and equality only for
a seven-edge spider matrix, hence the named tree. Conversely the
spider has three contributing mixed maximal independent sets.
\end{proof}

\paragraph{An essential non-exhaustion warning.}
The spider has nine size-four maximal independent sets.
They are not exhausted by the two pure classes and three mixed
$2+2$ sets. With the canonical bipartition, the sets
$\{1,2,3,7\}$ and $\{2,4,6,7\}$ are respectively $3+1$ and
$1+3$ maximal independent four-sets. They do not contribute to
$h(G)$, and their existence does not invalidate the equality
$h(G)=b(A)$.


\subsection{Why total MIS counts do not determine the statistic}
Liu's paper \cite{Liu1993} addresses the total maximal-independent-set
count; only its publisher abstract was accessible for this audit.
A source-level comparison is possible with Li--Zhang--Zhang
\cite[Lemma 3]{Li2013}, which records the classical extremal tree
families. Their family $B(2,3)$ consists of a core edge with three
length-two paths attached at its endpoints. Putting all three at one
endpoint gives our spider; putting one at one endpoint and two at the
other gives a different tree $T$.

For an explicit labelled $T$, take edges
$01,02,14,16,23,45,67$. Its maximal independent sets, grouped by
cardinality, have counts $3$ of size three and $6$ of size four.
The three size-three sets are $246,256,247$, and the size-four sets
are $0346,0356,0347,1257,0357,1357$ (digits denote vertices).
The only compatible triple is $\{0346,1257,0357\}$.
The independent finite checker confirms this list and its $2520$
labelling outcomes. The spider
and $T$ both have nine maximal independent sets, but
\[
\begin{array}{c|ccc}
 & |\mathrm{MIS}| & |\cF| & h\\\hline
S(2,2,2,1)&9&9&3\\
T&9&6&1
\end{array}
\]
Consequently their three-witness probabilities are respectively
$1/105$ and $1/315$. This proves that the total MIS count alone cannot
recover our statistic. It does not establish priority: the relevant
size-refined counts, mixed-set enumeration and design-incidence laws
still require comparison with the literature.

\section{The complete random cap-labelling law}
Fix an admissible eight-cap $U$, let $\lambda:V(G)\to U$ be a uniformly
chosen bijection, and write
\[
N_0=\#\left\{S\in\cF(G):\sum_{v\in S}\lambda(v)=0\right\}.
\]
Write $m=|\cF(G)|$, and for $t=0,2$ write $n_t$ for the
number of unordered pairs of distinct members of $\cF(G)$
whose intersection has cardinality $t$. Define
\[
M_1=\frac{2m}{35},\qquad
M_2=\frac{2n_0}{35}+\frac{n_2}{315},\qquad
M_3=\frac{h(G)}{315}.
\]

\Needspace{13\baselineskip}
\begin{theorem}[Full witness-count law]
For connected $G$, the random variable $N_0$ takes values in
$\{0,1,2,3\}$ and
\[
\begin{aligned}
\Pr(N_0=0)&=1-M_1+M_2-M_3,\\
\Pr(N_0=1)&=M_1-2M_2+3M_3,\\
\Pr(N_0=2)&=M_2-3M_3,\\
\Pr(N_0=3)&=M_3\le\frac1{105}.
\end{aligned}
\]
The last inequality is sharp exactly at $S(2,2,2,1)$.
\end{theorem}
\begin{proof}
The $8!$ bijections are equiprobable. A fixed four-set has each
of the $70$ image four-sets under exactly $4!4!$ bijections, and
four images are zero quartets. This gives $4/70=2/35$.
A complementary ordered pair has $70$ possible images, four favourable,
so has the same probability. An overlap-two ordered pair has
$70\binom42^2=2520$ images and eight favourable images. Each image
has exactly $(2!)^4=16$ preimage bijections, since the four membership
atoms have size two. Odd-overlap pairs cannot be two distinct zero
quartets. For a prescribed
compatible third set, there are
$\binom84\binom42^2=2520$ ordered abstract triples of the
relevant shape, with eight zero-quartet image triples. Again the
four membership atoms have size two, so every image triple has exactly
$16$ preimage bijections. Thus this is an actual uniform push-forward
of the bijection model, and the triple probability is $8/2520=1/315$.
Thus $\mathbb E\binom{N_0}i=M_i$ for $i=1,2,3$.
If all four quartets $A_i\cup B_j$ were independent, no edge
could join $A_1\cup A_2$ to $B_1\cup B_2$, because every such pair
occurs together in a quartet. Those are two nonempty complementary
parts of the vertex set, contradicting connectedness. Thus $N_0\le3$. Binomial moment inversion yields the displayed
probabilities. The extremal probability follows from $h\le3$.
\end{proof}

\section{Optional greedy dynamics and signed pebbling residues}
For any graph with at least one edge, independently run the
random-order greedy maximal-independent-set algorithm and sample a
uniform admissible-cap bijection. By symmetry, its zero-charge
probability is $\frac2{35}\Pr(|I|=4)$.
For rooted trees, three channel-wise residues
\[
Q_j=\sum_v(-1)^{d(r,v)}C_j(v)\pmod 3
\]
are conserved by the $2\to1$ adjacent-vertex pebbling move.
These supporting transfers are not claimed to imply a
reachability/stackability result.

\section{Formal verification and publication gate}
The three arbitrary Boolean-matrix extremal transfers and the concrete
spider witnesses were strictly kernel-verified at S003R8 commit
\href{https://github.com/jfairfaxball-348/Pan-Disciplinary-Invariance-Conjecture/commit/1e4cadd580b385e0d8e69e836f3211129e705839}{\texttt{1e4cadd580b3}}, hosted workflow
38887318919. All proofs use only standard Lean axioms. The S003R9 graph and connected-partition extensions passed at
\texttt{8199e6bc41dc}, workflow 38889327092; the exact certificate appears in the
\href{https://github.com/jfairfaxball-348/Pan-Disciplinary-Invariance-Conjecture/blob/main/sessions/S003R9_CLOSEOUT.md}{S003R9 closeout}.

\paragraph{Precisely scoped formal appendix.}
For each Boolean $4\times4$ matrix $A$, define a simple graph on
$[4]\sqcup[4]$ by making $i$ adjacent to $j$ across the sides exactly
when $A_{ij}=1$. Independence plus domination is proved equivalent to
maximality under inclusion. Disjoint union of the side sets yields a
bijection between mixed maximal independent vertex sets and eligible
left pairs. Consequently, for the actual vertex-set count $B(A)$,
\[
e(A)\ge6\Longrightarrow B(A)\le3,\qquad
[e(A)\ge7]\Longrightarrow
(B(A)=3\Longleftrightarrow A\text{ has the certified spider pattern}).
\]
The edge adapter identifies the actual unordered graph edge count with
$e(A)$, and connectedness supplies $e(A)\ge7$. Thus the connected
matrix graph has at most three mixed vertex-set MIS, with equality
exactly at the certified pattern. A separate
general graph argument proves uniqueness of proper Boolean colourings,
and hence complementary independent partitions, up to exchanging sides
in a connected graph. All these S003R9 statements passed strict hosted Lean workflow
38889327092 at revision \texttt{8199e6bc41dc}, using only standard axioms.

\paragraph{Remaining formal interface.}
The missing end-to-end links are: transport an arbitrary balanced
bipartition to the displayed matrix graph; identify its unordered
compatible-triple count with $B(A)$; construct the graph isomorphism from a spider-pattern matrix; and
realise the uniform $2520$-assignment model with its regular labelling
fibres and moments. The classical cap equivalence remains external.
The written theorems above do not depend on pretending these kernel
interfaces are already complete. This draft is not ready for submission:
priority, source comparison and final mathematical editing remain open.

\begin{thebibliography}{9}
\bibitem{Follett2014}
M.~Follett, K.~Kalail, E.~McMahon, C.~Pelland and R.~Won,
\emph{Partitions of $AG(4,3)$ into maximal caps},
Discrete Mathematics \textbf{337} (2014), 1--8.
\href{https://doi.org/10.1016/j.disc.2014.08.002}{doi:10.1016/j.disc.2014.08.002}.
\bibitem{Liu1993}
J.~Liu, \emph{Maximal independent sets in bipartite graphs},
Journal of Graph Theory \textbf{17} (1993), 495--507.
\href{https://doi.org/10.1002/jgt.3190170407}{doi:10.1002/jgt.3190170407}.
\bibitem{Li2013}
S.~Li, H.~Zhang and X.~Zhang,
\emph{Maximal independent sets in bipartite graphs with at least one cycle},
Discrete Mathematics \& Theoretical Computer Science \textbf{15}(2)
(2013), 243--258.
\href{https://dmtcs.episciences.org/607/pdf}{Open article}.
\end{thebibliography}
\end{document}
